\documentclass[tikz]{standalone}
\usepackage[utf8]{inputenc}
\usepackage{zxjatype}
\usepackage{xeCJK}
\usepackage{fontspec}
\usepackage{tikz}
\usetikzlibrary{shapes,arrows,positioning,calc}

% フォント設定（メインファイルと同じ）
\setCJKmainfont[
    BoldFont=BIZUDPGothic-Bold.ttf,
    Path=C:/Users/yutok/Downloads/BIZ_UDPGothic/
]{BIZUDPGothic-Regular.ttf}

\setmainfont[
    Path=C:/Users/yutok/Downloads/times-new-roman_freefontdownload_org/
]{times-new-roman.ttf}

\begin{document}
\begin{tikzpicture}[
    node distance=0.3cm,
    auto,
    thick,
    actor/.style={rectangle, draw=blue!70!black, fill=blue!5, text width=1.8cm, text centered, minimum height=0.7cm, font=\footnotesize, rounded corners},
    object/.style={rectangle, draw=gray!70!black, fill=gray!5, text width=1.8cm, text centered, minimum height=0.7cm, font=\footnotesize, rounded corners},
    arrow/.style={->, >=stealth, semithick, color=blue!60!black},
    return arrow/.style={->, >=stealth, semithick, dashed, color=gray!60!black},
    lifeline/.style={dashed, semithick, color=gray!40}
]

% 参加者
\node[actor] (user) at (0,0) {ユーザー};
\node[object] (app) at (3,0) {Flask App};
\node[object] (security) at (6,0) {Security\\Validator};
\node[object] (nlu) at (9,0) {NLU\\Processor};
\node[object] (scoring) at (12,0) {Scoring\\Engine};
\node[object] (db) at (15,0) {Database};

% ライフライン
\draw[lifeline] (user) -- (0,-6.5);
\draw[lifeline] (app) -- (3,-6.5);
\draw[lifeline] (security) -- (6,-6.5);
\draw[lifeline] (nlu) -- (9,-6.5);
\draw[lifeline] (scoring) -- (12,-6.5);
\draw[lifeline] (db) -- (15,-6.5);

% メッセージ（時系列順）
% 1. POSTリクエスト
\draw[arrow] (0,-0.8) -- (3,-0.8) node[midway, above, font=\tiny, sloped] {POST /};

% 2. セキュリティ検証
\draw[arrow] (3,-1.5) -- (6,-1.5) node[midway, above, font=\tiny, sloped] {validate\_input()};
\draw[return arrow] (6,-2.0) -- (3,-2.0) node[midway, below, font=\tiny, sloped] {is\_safe};

% 3. NLU処理
\draw[arrow] (3,-2.7) -- (9,-2.7) node[midway, above, font=\tiny, sloped] {extract\_symptoms()};
\draw[return arrow] (9,-3.2) -- (3,-3.2) node[midway, below, font=\tiny, sloped] {symptoms};

% 4. スコアリング
\draw[arrow] (3,-3.9) -- (12,-3.9) node[midway, above, font=\tiny, sloped] {calculate\_score()};
\draw[arrow] (12,-4.4) -- (15,-4.4) node[midway, above, font=\tiny, sloped] {get\_medicines()};
\draw[return arrow] (15,-4.9) -- (12,-4.9) node[midway, below, font=\tiny, sloped] {medicines};
\draw[return arrow] (12,-5.4) -- (3,-5.4) node[midway, below, font=\tiny, sloped] {scores};

% 5. 最終結果
\draw[return arrow] (3,-6.0) -- (0,-6.0) node[midway, below, font=\tiny, sloped] {recommendation};

\end{tikzpicture}
\end{document}
